\documentclass{article}
\usepackage{graphicx}
\usepackage{amsmath}
\usepackage[T2A]{fontenc}  % For Cyrillic fonts
\usepackage[utf8]{inputenc}  % For UTF-8 encoding
\usepackage[russian,english]{babel}  % For Russian language support
\usepackage{amssymb}
\usepackage[margin=2cm]{geometry}
\usepackage{listings}
\usepackage{xcolor}
\usepackage{float}
\usepackage{hyperref}

\lstset{
	language=Python,
	basicstyle=\ttfamily\small,
	keywordstyle=\color{blue},
	stringstyle=\color{red},
	commentstyle=\color{gray},
	backgroundcolor=\color{lightgray!20},
	frame=single,
	rulecolor=\color{black},
	numbers=left,
	numberstyle=\tiny\color{gray},
	stepnumber=1,
	numbersep=10pt,
	showspaces=false,
	showstringspaces=false,
	breaklines=true,
	breakatwhitespace=true,
	tabsize=4,
	captionpos=b
}

\title{Математическая статистика \\ Домашнее задание 3}
\author{Никита Загайнов}

\date{Май 2025}

\begin{document}

\maketitle

\begin{enumerate}
	\item \begin{enumerate}
		      \item
		            Воспользуемся критерием Пирсона для проверки гипотезы о том, что вероятность прибытия пациента в любой из дней недели равномерна. Постановка теста:

		            \begin{equation*}
			            \begin{aligned}
				            H_0 & : p_i = \frac{1}{7} \text{, } \forall i \in [1,7]    \\
				            H_1 & : p_i \neq \frac{1}{7} \text{, } \exists i \in [1,7]
			            \end{aligned}
		            \end{equation*}

		            Статистика теста:
		            \[
			            \chi^2 = n \sum_{i=1}^{7} \frac{\left( \frac{\nu_i}{n} - p_i \right)^2}{p_i}
		            \]
		            где $n$ - общее количество наблюдений, $\nu_i$ - количество наблюдений в $i$-й категории. Если нулевая гипотеза верна, то $\chi^2 \sim \chi^2_6$, в противном случае $\chi^2$ стремится к бесконечности с ростом $n$. Тогда тест будет проводиться следующим образом:
		            \begin{equation*}
			            \begin{aligned}
				            \chi^2 \leq \chi^2_6(0.95) & \Rightarrow \text{отвергаем } H_1 \\
				            \chi^2 > \chi^2_6(0.95)    & \Rightarrow \text{отвергаем } H_0
			            \end{aligned}
		            \end{equation*}

		            \[
			            \chi^2 \approx 15.389 > \chi^2_6(0.95) \approx 12.592 \Longrightarrow \text{отвергаем } H_0
		            \]

		      \item Код решения:
		            \begin{lstlisting}
import numpy as np
from scipy import stats


nu = np.array([36, 53, 35, 26, 30, 44, 28])
n = np.sum(nu)
m = len(nu)
p0 = np.ones(m) / m
alpha = 0.05

statistic = n * np.sum((nu / n - p0) ** 2 / p0)

print(f"Criterion statistic: {statistic:.3f}")

quantile = stats.chi2.ppf(1 - alpha, m - 1)

print(f"Quantile of Chi squared distribution: {quantile:.3f}")

if statistic > quantile:
    print("Reject the null hypothesis")
else:
    print("Reject the alternative hypothesis")
                    \end{lstlisting}

		            Вывод:
		            \begin{lstlisting}
Criterion statistic: 15.389
Quantile of Chi squared distribution: 12.592
Reject the null hypothesis
                    \end{lstlisting}
	      \end{enumerate}

	\item
	      Найдём ОМП для параметра $\theta$. Функция правдоподобия:
	      \[
		      L(\theta) = ((1 - \theta)^2)^{\frac{n}{4}} (2\theta(1 - \theta))^{\frac{n}{4}} (\theta^2)^{\frac{n}{2}}
	      \]
	      Возьмём логарифм:
	      \[
		      \log L(\theta) = \frac{n}{2} \log(1 - \theta) + \frac{n}{4} \log(2\theta(1 - \theta)) + n \log(\theta)
	      \]
	      Возьмём производную и приравняем к нулю:
	      \[
		      \frac{d \log L(\theta)}{d \theta} = -\frac{n}{2(1 - \theta)} + \frac{n(1 - 2\theta)}{4\theta(1 - \theta)} + \frac{n}{\theta} = 0
	      \]
	      Решив уравнение, получаем:
	      \[
		      \theta = \frac{5}{8}
	      \]

	      Теперь проверим гипотезу о том, что выборка следует полученному распределению:
	      \begin{equation*}
		      \begin{aligned}
			      H_0 & : p = \begin{bmatrix}
				                  \left(\frac{3}{8}\right)^2        \\[5pt]
				                  2\cdot\frac{3}{8}\cdot\frac{5}{8} \\[5pt]
				                  \left(\frac{5}{8}\right)^2
			                  \end{bmatrix} = \begin{bmatrix}
				                                  \frac{9}{64}  \\[5pt]
				                                  \frac{30}{64} \\[5pt]
				                                  \frac{25}{64}
			                                  \end{bmatrix} \\
			      H_1 & : p \neq \begin{bmatrix}
				                     \frac{9}{64}  \\[5pt]
				                     \frac{30}{64} \\[5pt]
				                     \frac{25}{64}
			                     \end{bmatrix}
		      \end{aligned}
	      \end{equation*}

	      Воспользуемся критерием Пирсона:
	      \[
		      \chi^2 = n \sum_{i=1}^{3} \frac{\left( \frac{\nu_i}{n} - p_i \right)^2}{p_i} \approx 27.876 > \chi^2_2(0.95) \approx 5.991 \Longrightarrow \text{отвергаем } H_0
	      \]

	\item
		Решение задачи находится в блокноте во вложении

\end{enumerate}


\end{document}
